\chapter{Descrizione stage}
\label{cap:descrizione-stage}

\intro{In questo capitolo vengono descritte le attività 
    svolte, l'organizzazione dello stage e l'analisi dei rischi.}\\

\section{Obiettivi del progetto}
Il tirocinio ha come obiettivo principale quello di avvicinare studentesse e 
studenti al mondo del lavoro, apprendendo le tecnologie e gli strumenti 
maggiormente richiesti in ambito aziendale.
Al termine del tirocinio svolto in RiskApp S.r.l. erano attese le seguenti competenze:
\begin{itemize}
    \item progettazione e implementazione di un server conforme al \gls{mcpg} per l'esposizione
    di funzionalità tramite \gls{toolg} dedicati;
    \item implementazione di tecniche di validazione e \gls{prompt-engineeringg} per 
    il controllo della qualità dei risultati prodotti dal \gls{llmg}\glsfirstoccur ;
    \item implementazione di un sistema di \gls{cachingg} intelligente per la gestione 
    efficiente delle richieste ridondanti;
    \item sviluppo di un sistema di autenticazione e autorizzazione sicuro 
    basato su \gls{jwtg}\glsfirstoccur ;
    \item utilizzo dello strumento di containerizzazione \gls{dockerg}\glsfirstoccur per la gestione 
    e l'orchestrazione dei servizi;
    \item gestione di database per la persistenza dei dati relativi alle richieste 
    e agli utenti;
    \item sviluppo di un'interfaccia utente intuitiva e reattiva.
\end{itemize}

\section{Vincoli}
L'azienda ha definito alcune linee guida per 
garantire la portabilità del prodotto finale e l'interoperabilità con i software 
aziendali interni. I vincoli individuati sono i seguenti:
\begin{itemize}
    \item il prodotto deve essere sviluppato usando il \gls{frameworkg}\glsfirstoccur  React 
    per il \textit{frontend};
    \item il prodotto deve essere sviluppato usando il \gls{frameworkg} FastAPI
    per il \textit{backend};
    \item il sistema deve essere containerizzato tramite \gls{dockerg};
    \item le componenti del sistema devono comunicare tramite \gls{apig}\glsfirstoccur .
\end{itemize}

\section{Pianificazione}
All'inizio del tirocinio è stato redatto il piano di lavoro per 
la ripartizione approssimativa delle attività su base settimanale.
La pianificazione prevista è dunque la seguente:
\begin{itemize}
    \item \textbf{Prima settimana} 
        \begin{itemize}
            \item incontro con le persone coinvolte nel progetto per discutere 
            i vari requisiti;
            \item studio e formazione sulle tecnologie adottate, 
            quali FastMCP e FastAPI.
        \end{itemize}
    \item \textbf{Seconda settimana }
        \begin{itemize}
            \item studio e formazione sulle tecnologie adottate;
            \item analisi dei requisiti.
        \end{itemize}
    \item \textbf{Terza settimana }
        \begin{itemize}
            \item sviluppo del prototipo;
            \item inizio progettazione;
            \item inizio pianificazione.
        \end{itemize}
    \item \textbf{Quarta settimana }
        \begin{itemize}
            \item configurazione dell'ambiente di sviluppo con \gls{dockerg};
            \item implementazione del componente server \gls{mcpg}.
        \end{itemize}
    \item \textbf{Quinta settimana }
        \begin{itemize}
            \item implementazione della conversione delle pagine in 
            formato \gls{markdowng} e \gls{htmlg};
            \item implementazione del sistema di \gls{cachingg} intelligente.
        \end{itemize}
    \item \textbf{Sesta settimana }
        \begin{itemize}
            \item gestione degli errori e casi limite;
            \item ottimizzazione delle performance del sistema.
        \end{itemize}
    \item \textbf{Settima settimana}
        \begin{itemize}
            \item sviluppo dei \textit{test};
            \item collaudo e verifica;
        \end{itemize}
    \item \textbf{Ottava settimana}
        \begin{itemize}
            \item completamento del collaudo e correzione \textit{bug};
            \item stesura della documentazione tecnica finale.
        \end{itemize}
\end{itemize}

\subsection{Prima settimana}
Durante la prima settimana di lavoro ha avuto luogo l'incontro iniziale con i 
\textit{tutor} aziendali per un'analisi approfondita del progetto e la definizione 
dei requisiti principali. Tale incontro è stato fondamentale per definire con precisione 
gli obiettivi e delineare i vincoli tecnologici e operativi. 
Parallelamente, è iniziata una fase di studio individuale focalizzata 
sulla lettura della documentazione relativa ai \gls{frameworkg} 
\textit{FastMCP} e \textit{FastAPI}. Per velocizzare l'apprendimento di tali 
tecnologie ho sviluppato dei piccoli progetti guidati in \textit{Python} 
disponibili nella documentazione ufficiale (\textit{QuickStart}), 
Tali progetti erano volti alla comprensione pratica delle tecnologie 
studiate. In particolare:
\begin{itemize}
    \item riguardo \textbf{FastAPI} ho imparato a implementare:
        \begin{itemize}
            \item funzioni \texttt{async}/\texttt{await};
            \item classi con \texttt{@dataclass};
            \item \textit{streaming} di dati in formato \gls{jsong} per il \gls{sseg}\glsfirstoccur ;
            \item \gls{websocketg};
            \item i principali metodi \gls{httpg} (\texttt{GET}, \texttt{POST}, 
            \texttt{PUT}, \texttt{DELETE});
        \end{itemize}
    \item riguardo \textbf{FastMCP}, ho imparato a implementare:
        \begin{itemize}
            \item \gls{toolg} \gls{mcpg};
            \item \textit{resources} \gls{mcpg};
            \item \textit{prompts} \gls{mcpg};
        \end{itemize}
\end{itemize}


\subsection{Seconda settimana}
La seconda settimana è stata dedicata al consolidamento delle competenze tecnologiche 
acquisite nei giorni precedenti. Per una comprensione più solida del progetto, si è 
proceduto con la redazione del documento di Analisi dei Requisiti. Tale attività ha 
permesso di formalizzare tutti i requisiti, che sono stati suddivisi in requisiti funzionali,
di qualità e di vincolo e a sua volta scomposti in obbligatori, opzionali e desiderabili.
Durante questo periodo sono 
stati svolti alcuni incontri con i \textit{tutor} aziendali per chiarire alcuni dubbi
e aggiungere nuovi requisiti, i quali hanno permesso di espandere il progetto.

\subsection{Terza settimana}
Nel corso della terza settimana l'attenzione si è spostata verso le attività di sviluppo 
di un \textit{Proof of Concept}, volto a permettere all'azienda di valutare la mia comprensione degli obiettivi
e dell'attuabilità del progetto. 
\\Il prototipo realizzato ha previsto lo sviluppo 
del server \gls{mcpg} che esponeva un 
\gls{toolg} di \gls{crawlingg}.
Tale \gls{toolg} richiedeva come parametri il nome del sito web da cui 
eseguire il \gls{crawlingg}, la domanda dell'utente in linguaggio naturale 
e il tipo di formato desiderato per la generazione della risposta.

Il sistema avviava un processo di \gls{crawlingg} che, ricorsivamente, 
estraeva l'\gls{htmlg} della pagina, lo riassumeva e tentava di rispondere 
alla domanda dell'utente.
Se le informazioni ricavate fino a quel punto non erano sufficienti a 
rispondere al quesito, il sistema prelevava tutti i \textit{link} 
ai siti del medesimo 
dominio web, selezionava i più pertinenti (al 
più tre) e proseguiva la ricerca.
Per evitare una ricorsione infinita o un numero eccessivo di 
chiamate ricorsive il sistema definiva due parametri: l'ampiezza massima (il 
numero massimo di \textit{link} selezionabili ad ogni livello di 
ricorsione) e profondità massima (numero 
massimo di livelli di ricorsione).
Successivamente è stata sviluppata anche una semplice interfaccia utente 
per l'inserimento dei parametri e la visualizzazione della risposta generata.
Infine, è stato implementato un processo di \textit{streaming} tramite \gls{sseg}, 
per mostrare nella pagina web i singoli passaggi svolti durante la ricorsione 
(estrazione della pagina, riassunto del contenuto, estrazione di tutti i \textit{link}, 
scelta dei \textit{link} più pertinenti).

Il prototipo sviluppato ha permesso di dimostrare concretamente la fattibilità 
dell'approccio proposto, fornendo all'azienda una prima base tangibile su cui valutare 
la direzione del progetto e indirizzare le fasi di sviluppo successive.

In base ai riscontri del \textit{Proof of Concept}, è iniziata la pianificazione 
di dettaglio dell'architettura del sistema, prima per il 
\textit{backend} e successivamente per il \textit{frontend}.
In particolare, sono state effettuate le seguenti scelte architetturali:
\begin{itemize}
    \item suddivisione del \textit{backend} in microservizi indipendenti, 
    ciascuno responsabile di un dominio applicativo specifico (autenticazione, 
    gestion utenti, storico delle richieste, ricerca e server \gls{mcpg}) con 
    l'obiettivo di favorire la scalabilità e la manutenibilità del sistema;
    \item adozione dell'architettura esagonale (\textit{hexagonal architecture})
    all'interno di ciascun microservizio al fine di disaccoppiare la logica di business
    dai dettagli implementativi (come database, \gls{llmg}, \gls{apig} esterne) e
    favorire la testabilità del codice;
    \item adozione del pattern \gls{mvvmg}\glsfirstoccur per il \textit{frontend} allo scopo di 
    separare la logica di presentazione dalla gestione dello stato dell'applicazione.
\end{itemize}


\subsection{Quarta settimana}
La quarta settimana ha segnato l'inizio della fase di codifica dell'infrastruttura, 
partendo dal componente \textit{backend}. Il primo obiettivo è stato la 
configurazione dell'ambiente di sviluppo attraverso la containerizzazione 
con \gls{dockerg} e la suddivisione dei container per microservizi, 
al fine di garantire fin dall'inizio la portabilità e 
scalabilità del software, uno dei vincoli principali imposti dall'azienda.
In seguito è stato definito la base di dati dell'applicazione, al fine di 
garantire la persistenza dei dati relativi agli utenti (profilo, permessi, 
\text{chat}) e alle richieste effettuate al \gls{llmg}.
\\Infine, è stata avviata la codifica dei vari microservizi 
del \textit{backend}. Il processo di codifica di tali microservizi è 
stato protratto per le successive settimane e sono stati 
sviluppati nel seguente ordine:
\begin{enumerate}
    \item \textbf{server \gls{mcpg} service}: rappresenta il vero e proprio servizio \gls{mcpg} che espone il \gls{toolg} di \textit{crawl} e per ogni passaggio effettuato,
    invia al \textit{frontend} un \textit{event} tramite \gls{sseg};
    \item \textbf{search service}: rappresenta il servizio dedicato alla gestione dei messaggi dell'utente il linguaggio naturale. Il microservizio, dato il contesto, comprende quale \gls{toolg} è richiesto dall'utente, estrae i parametri necessari e chiama l'\textit{endpoint} corretto del microservizio Server MCP;
    \item \textbf{auth service}: rappresenta il servizio dedicato all'autenticazione dell'utente e alla gestione dei \gls{jwtg} (validazione e creazione di \textit{access token} e \textit{refresh token}, salvataggio e revoca del \textit{refresh token});
    \item \textbf{user service}: rappresenta il servizio dedicato alla creazione e 
    alla gestione degli utenti;
    \item \textbf{history service}: rappresenta il servizio dedicato alla gestione 
    e salvataggio dei messaggi delle \textit{chat} tra l'utente e l'\gls{llmg}.
\end{enumerate}
Dalla quarta settimana la pianificazione prevista dal piano di lavoro non è stata più 
seguita: dopo l'analisi dei requisiti sono sorte nuove funzionalità importanti e 
complesse che non erano state dichiarate all'inizio. Il prodotto finale avrebbe 
dunque richiesto più lavoro su vari fronti, diversi dalla sola gestione di un servizio 
\gls{mcpg}. Durante tale periodo è stato dunque fondamentale una ripianificazione 
completa delle settimane a seguire in maniera tale da adattarla alle nuove specifiche.

\subsection{Quinta settimana}
Durante la quinta settimana si è continuato il processo di codifica \textit{backend}avviato nella quarta settimana. In particolare:
\begin{itemize}
    \item è stato completato il microservizio relativo al server \gls{mcpg};
    \item è stato avviato e completato \textit{search service};
    \item è stato avviato \textit{auth service}.
\end{itemize}

\subsection{Sesta settimana}
Nella sesta settimana si è continuato il processo di codifica, in particolare:
\begin{itemize}
    \item è stato completato il microservizio \textit{auth service};
    \item sono stato avviati e completati \textit{user service} 
    e \textit{history service}.
\end{itemize}
Avendo concluso lo sviluppo del \textit{backend}, ho iniziato 
la codifica relativa alla parte del \textit{frontend}, la quale è stata 
suddivisa in base alle funzionalità fornite. Sono state dunque 
implementate le seguenti pagine web:
\begin{itemize}
    \item \textbf{login} per effettuare l'accesso al sistema tramite le 
    credenziali \textit{username} e \textit{password};
    \item \textbf{chatbot} per inviare messaggi al sistema e visualizzare 
    la risposta generata con i vari passaggi del processo di \gls{crawlingg};
    \item \textbf{gestione utenti}, dedicato ai soli utenti \textit{admin}, 
    permette di creare e gestire nuovi utenti all'interno del sistema;
    \item \textbf{profilo utente} per visualizzare e modificare le 
    informazioni dell'utente;
    \item \textbf{storico delle \textit{chat}} per visualizzare le varie 
    conversazioni passate.
\end{itemize}

\subsection{Settima settimana}
La settima settimana è stata dedicata al \textit{refactoring}: 
partendo dal \textit{backend} fino a raggiungere il \textit{frontend}, sono state analizzate tutte le componenti e le funzioni implementate e 
sono state ottimizzate, al fine di renderle 
più efficienti dal punto di vista di complessità, velocità e utilizzo di 
risorse.
Un'attività di particolare rilievo ha riguardato l'implementazione del sistema di \textit{tracing} all'interno del servizio \gls{mcpg}, realizzato mediante la piattaforma \textit{LangSmith}. Tale funzionalità consente di monitorare e analizzare in modo dettagliato le chiamate effettuate ai modelli di intelligenza artificiale impiegati dal sistema.


Sono state poi implementate piccole \textit{feature} aggiuntive per 
migliorare l'\textit{user experience}, come l'utilizzo di icone nei 
bottoni, il tema chiaro/scuro e la possibilità di copiare o 
scaricare la risposta generata. 
Infine sono stati sviluppati i vari \textit{test} relativi a tutte le componenti 
del sistema per validare il loro corretto funzionamento e valutare 
la conformità rispetto ai requisiti definiti in fase di analisi.

\subsection{Ottava settimana}
Durante l'ultima settimana sono stati corretti alcuni \textit{bug} 
relativi alla visualizzazione dei messaggi nelle \textit{chat} e
si è concluso ufficialmente il progetto con il collaudo finale 
da parte dei \textit{tutor} aziendali.


\subsection{Pianificazione in retrospettiva}
L'esperienza di tirocinio si è rivelata particolarmente formativa soprattutto dal 
punto di vista tecnico. Sebbene la pianificazione iniziale
abbia costituito un utile punto di partenza per organizzare le attività, il progetto
ha subito un'evoluzione importante durante il suo sviluppo. Dopo la fase di
analisi dei requisiti, avviata nella seconda settimana, sono infatti emerse nuove esigenze e funzionalità che hanno portato a una revisione significativa degli obiettivi iniziali e della pianificazione prevista.

Questa situazione ha evidenziato come, all'interno di un contesto aziendale reale,
i requisiti possano evolvere nel tempo in funzione delle necessità del cliente,
delle opportunità individuate durante lo sviluppo e dei risultati ottenuti nelle
fasi preliminari del progetto. Di conseguenza, è stato necessario adattare
costantemente il piano di lavoro, ridefinendo priorità e attività da svolgere.
Tale esperienza ha permesso di comprendere l'importanza della flessibilità e della
capacità di gestione del cambiamento, competenze fondamentali nello sviluppo
software moderno.

Dal punto di vista tecnico, il progetto ha rappresentato un'importante occasione
di approfondimento di tecnologie e paradigmi architetturali non affrontati in modo
approfondito durante il percorso universitario. In particolare, ho acquisito
competenze nell'utilizzo di Python e dei relativi \gls{frameworkg} FastAPI e FastMCP, nella progettazione di architetture a microservizi, nell'applicazione dell'architettura esagonale e nell'utilizzo del pattern MVVM per lo sviluppo \textit{frontend}. Ho inoltre maturato esperienza nella containerizzazione delle applicazioni mediante \gls{dockerg} e nella progettazione di sistemi che integrano modelli linguistici di grandi dimensioni (\gls{llmg}).

Un altro aspetto particolarmente significativo è stato il confronto continuo con
i \textit{tutor} aziendali. Le revisioni periodiche e le discussioni sulle scelte
progettuali hanno contribuito a migliorare la qualità delle soluzioni sviluppate
e mi hanno permesso di imparare molte tecnologie e tecniche nuove. 
Questo confronto costante ha inoltre favorito lo sviluppo
delle capacità comunicative.

Il raggiungimento degli obiettivi finali e la realizzazione di un prodotto completo,
composto da un'infrastruttura \textit{backend} a microservizi e da un'interfaccia
\textit{frontend} dedicata, hanno rappresentato una conferma concreta delle
competenze acquisite durante il percorso universitario e di quelle sviluppate nel
corso del tirocinio. L'esperienza ha quindi contribuito in maniera significativa
alla mia crescita professionale, fornendomi una visione più completa delle attività
necessarie per progettare, sviluppare, testare e distribuire un'applicazione software
all'interno di un contesto aziendale reale.


\section{Organizzazione del lavoro}
\subsection{Metodologia adottata}
Per la gestione e l'esecuzione delle attività durante tutto il periodo di stage è stata scelta una metodologia di sviluppo di tipo \textit{iterativo-incrementale} basata su \textit{timebox} settimanali. Questo approccio ha permesso di tradurre la pianificazione strategica di alto livello in obiettivi operativi intermedi e costantemente verificabili. La misura dell'avanzamento del progetto è stata determinata dal progressivo rilascio di incrementi software stabili e funzionanti, garantendo al contempo la flessibilità necessaria per ripianificare gli obiettivi in base alle esigenze emerse in corso d'opera.

Il modello iterativo-incrementale si fonda sul principio del miglioramento continuo del prodotto attraverso cicli ripetuti di analisi, progettazione, implementazione e verifica. Differisce dai modelli sequenziali poiché l'applicazione non viene sviluppata in un'unica soluzione finale, ma cresce per ''incrementi'' successivi. Ogni iterazione produce una versione del software che aggiunge nuove funzionalità o ottimizza quelle esistenti, consentendo di integrare tempestivamente i \textit{feedback} degli \textit{stakeholder} e mitigare i rischi tecnici.

Il ciclo di vita del progetto di tirocinio è stato mappato sulle quattro macro-fasi logiche caratteristiche di questo modello di sviluppo:
\begin{itemize}
\item \textbf{Inception}: svoltasi prevalentemente durante la prima settimana, ha compreso la definizione del perimetro del progetto, l'analisi preliminare dei vincoli tecnologici (\textit{Docker}, \textit{Python} con \textit{FastAPI} e \textit{FastMCP}, \textit{React}, \textit{PostgreSQL}) e l'allineamento iniziale con i \textit{tutor} aziendali per delineare gli obiettivi ad alto livello;
\item \textbf{Elaboration}: coincidente con la seconda e la terza settimana, in cui si è proceduto alla formalizzazione del documento di Analisi dei Requisiti e alla successiva progettazione architetturale, validata attraverso la realizzazione di un \textit{Proof of Concept};
\item \textbf{Construction}: estesa dalla quarta alla settima settimana, ha rappresentato la fase a più alta intensità di sviluppo. In questo periodo l'architettura a microservizi del \textit{backend} e l'interfaccia \textit{frontend} in \textit{React} sono state progressivamente codificate, integrate e containerizzate tramite \textit{Docker}. Sono stati sviluppati i \textit{test} automatici ed è stata svolta l'attività di \textit{refactoring} di tutte le componenti di sistema;
\item \textbf{Transition}: focalizzata nell'ultima settimana, ha riguardato l'ottimizzazione delle performance e la consegna finale del prodotto ai \textit{tutor} aziendali.
\end{itemize}

\subsection{Analisi comparativa e motivazioni della scelta}
\subsubsection{Gestione del cambiamento}
La scelta di questa metodologia è stata determinata dalla natura innovativa e dinamica del progetto, incentrato sull'utilizzo del protocollo \gls{mcpg} e sull'integrazione di modelli linguistici di grandi dimensioni (\gls{llmg}).

L'efficacia della scelta iterativa si è manifestata concretamente a partire dalla quarta settimana: l'evoluzione dei requisiti ha reso obsoleta la pianificazione iniziale. In un \gls{frameworkg} rigido, tale variazione avrebbe comportato una seria minaccia per la stabilità e le tempistiche del progetto; l'approccio iterativo ha invece permesso di avviare una ripianificazione delle settimane rimanenti, rimodulando i confini dei \textit{timebox} e garantendo lo sviluppo delle nuove funzionalità senza compromettere la qualità del software o i tempi di consegna.

\subsubsection{Confronto con altre metodologie}
Per valutare la qualità nel contesto del progetto, la metodologia adottata è stata messa a confronto con i principali paradigmi alternativi:
\begin{itemize}
\item \textbf{Modello Waterfall (A Cascata)}: questo approccio prevede una sequenza rigida di fasi in cui i requisiti devono essere interamente definiti e congelati all'inizio. Applicare il modello a cascata in questo scenario sarebbe stato fallimentare, poiché l'evoluzione dei requisiti avrebbe richiesto un ritorno alle fasi iniziali, provocando ritardi insostenibili.
\item \textbf{Scrum puro}: sebbene Scrum sia un'ottima metodologia agile per gestire requisiti fluttuanti, la sua struttura cerimoniale rigorosa (comprendente \textit{Daily Standup}, \textit{Sprint Planning}, \textit{Review} e \textit{Retrospective}) unita alla definizione formale di ruoli rigidi (\textit{Scrum Master}, \textit{Product Owner}) avrebbe generato un sovraccarico organizzativo ingiustificato per un progetto svolto individualmente.
\item \textbf{Kanban}: pur offrendo flessibilità e concentrandosi sul flusso continuo di lavoro, Kanban non impone scadenze temporali rigide o cicli di pianificazione predefiniti. L'assenza di \textit{timebox} strutturati avrebbe reso complessa la verifica periodica degli obiettivi soddisfatti.
\end{itemize}

\subsection{Metodologia operativa e flussi di comunicazione}
Le attività settimanali sono state governate da due dinamiche principali di sincronizzazione con la struttura aziendale:
\begin{enumerate}
    \item \textbf{Supporto quotidiano e comunicazione asincrona}: anche durante le giornale svolte da remoto, ero sempre in contatto con i \textit{tutor} aziendali tramite l'applicazione di messaggistica \textit{Slack}. 
    Tale canale ha permesso di mantenere un flusso informativo costante con i \textit{tutor}, facilitando la rapida risoluzione di problemi o di domande relative al progetto;
    \item \textbf{Allineamento finale settimanale}: un incontro della durata di circa 20 minuti programmato al termine di ogni ciclo. Durante la sessione veniva effettuato un controllo dei risultati ottenuti durante la settimana e veniva pianificato l'obiettivo del ciclo successivo.
\end{enumerate}

\subsection{Esempio di iterazione settimanale}
A titolo esemplificativo, nella Tabella~\ref{tab:esempio-iterazione} viene riportata la scheda di gestione utilizzata per documentare e tracciare le attività della terza settimana, periodo in cui è stato sviluppato il prototipo del sistema di \gls{crawlingg}.
\begin{table}[htbp]
    \centering
    \small
    \renewcommand{\arraystretch}{1.4}
    \begin{tabular}{|p{3cm}|p{10cm}|}
        \hline
        \textbf{Elemento} & \textbf{Descrizione} \\ \hline

        \textbf{Settimana} & Settimana 3 \\ \hline

        \textbf{Goal} & Sviluppo di un \textit{Proof of Concept} (PoC) funzionante, finalizzato alla verifica della fattibilità del server MCP. \\ \hline

        \textbf{Attività pianificate in ordine di priorità} &
        \begin{enumerate}
            \item Implementazione base del server \textit{MCP} con esposizione di un \textit{tool} di \textit{crawling};
            \item Sviluppo della logica di estrazione, riassunto e selezione \textit{link} delle pagine HTML;
            \item Formattazione della risposta generata nei formati HTML, JSON, plain text e Markdown con vincoli di ampiezza e profondità massima;
            \item Realizzazione di un’interfaccia utente di base con inserimento URL dei sito web, inserimento domanda e visualizzazione della risposta finale;
            \item Implementazione dello \textit{streaming} degli stati tramite \textit{Server-Sent Events} (SSE);
            \item Visualizzazione via \textit{frontend} dei passaggi in tempo reale.
        \end{enumerate}
        \\ \hline

        \textbf{Esito} &
        Obiettivo raggiunto. \\ \hline

        \textbf{Feedback} &
        Valutazione positiva. Viene consigliato di rendere \textit{built-in} (statici) i parametri di massima ampiezza e profondità. \\ \hline

    \end{tabular}
    \caption{Documentazione di un’iterazione settimanale (Settimana 3).}
    \label{tab:esempio-iterazione}
\end{table}


\section{Analisi preventiva dei rischi}

Durante la fase di analisi iniziale sono stati individuati alcuni possibili rischi a cui si potrà andare incontro.
Si è quindi proceduto a elaborare delle possibili soluzioni per far fronte a tali rischi.\\

\begin{risk}{Elevata curva di apprendimento tecnologico}
    \riskdescription{L'adozione di un ampio \textit{stack} tecnologico comprendente \gls{frameworkg} e paradigmi non approfonditi nel percorso accademico (FastAPI, protocollo \gls{mcpg}, architettura esagonale e \gls{dockerg}) rischia di sottrarre tempo prezioso alla fase più dispendiosa, ovvero la codifica}
    \risksolution{Pianificazione di una fase iniziale di formazione teorico-pratica con sviluppo di progetti guidati (\textit{QuickStart}) e di un \textit{Proof of Concept} per validare le competenze prima dello sviluppo definitivo}
    \label{risk:curva-apprendimento} 
\end{risk}

\begin{risk}{Volatilità e instabilità dei requisiti}
    \riskdescription{Trattandosi di un'applicazione legata a tecnologie emergenti in un contesto aziendale reale, i requisiti iniziali potrebbero subire variazioni o ampliamenti in corso d'opera, invalidando la pianificazione strategica originaria}
    \risksolution{Adozione di una metodologia iterativo-incrementale basata su \textit{timebox} settimanali e canali di comunicazione costanti (\textit{Slack} e allineamenti), permettendo una tempestiva ripianificazione degli obiettivi}
    \label{risk:volatilita-requisiti} 
\end{risk}

\begin{risk}{Latenza e limitazioni nell'integrazione di LLM e crawling}
    \riskdescription{L'interazione con i modelli linguistici e i processi ricorsivi di \textit{web crawling} potrebbero introdurre colli di bottiglia e problemi di \textit{rate limiting} delle \textit{API}}
    \risksolution{Definizione di parametri rigidi di ampiezza e profondità massima per la ricorsione, implementazione di un sistema di \textit{caching} intelligente e utilizzo della piattaforma \textit{LangSmith} per il tracciamento e il monitoraggio continuo delle chiamate}
    \label{risk:latenza-llm} 
\end{risk}

\begin{risk}{Non-determinismo delle risposte del modello linguistico (LLM)}
    \riskdescription{I modelli linguistici di grandi dimensioni possono generare risposte imprevedibili o strutturalmente non omogenee (comunemente chiamate ''allucinazioni'') a parità di \textit{input}, rischiando di compromettere la consistenza dei dati inviati al \textit{frontend} o la corretta esecuzione dei \gls{toolg} basati sul protocollo \gls{mcpg}}
    \risksolution{Implementazione di tecniche stringenti di \textit{prompt engineering}, definizione di vincoli sintattici rigorosi sull'\textit{output} del modello (richiedendo formati strutturati, come \gls{jsong}) e aggiunta di validazione dei dati per intercettare e gestire tempestivamente eventuali risposte non conformi}
    \label{risk:non-determinismo-llm} 
\end{risk}

\begin{risk}{Complessità di integrazione dell'architettura a microservizi}
    \riskdescription{La scomposizione del \textit{backend} in microservizi indipendenti introduce complessità nella comunicazione tra componenti e nella gestione della persistenza}
    \risksolution{Configurazione e orchestrazione dell'ambiente di sviluppo tramite \textit{Docker} fin dalle prime fasi e sviluppo sequenziale e incrementale dei singoli microservizi}
    \label{risk:integrazione-microservizi} 
\end{risk}


\section{Requisiti e obiettivi}
\label{sec:requisiti-obiettivi}

In questa sezione vengono formalizzati gli obiettivi prefissati all'inizio dell'esperienza di tirocinio, classificati in base alla loro priorità di realizzazione.

\subsection{Obiettivi obbligatori}
Gli obiettivi obbligatori rappresentano i requisiti minimi e indispensabili per il successo del progetto. Essi comprendono le componenti architetturali e le funzionalità primarie di interazione con i modelli linguistici:

\begin{itemize}
    \item \textbf{Progettazione e implementazione del server MCP}: Sviluppo di un server conforme allo standard \textit{Model Context Protocol} (MCP) mediante l'utilizzo del framework \textit{FastMCP}, finalizzato all'esposizione di funzionalità e strumenti dedicati all'ecosistema di intelligenza artificiale.
    \item \textbf{Sviluppo del modulo di crawling parametrico}: Realizzazione di un motore di \textit{web crawling} ricorsivo in grado di estrarre e riassumere il contenuto informativo di pagine web in linguaggio naturale, governato da vincoli rigidi di ampiezza e profondità massima per prevenire cicli infiniti o sovraccarichi computazionali.
    \item \textbf{Ingegnerizzazione dei prompt e validazione dell'output AI}: Applicazione di tecniche avanzate di \textit{prompt engineering} e introduzione di uno strato software di validazione sintattica e semantica, volto a garantire la qualità e la coerenza dei risultati prodotti dal \textit{Large Language Model} (LLM).
    \item \textbf{Architettura del backend a microservizi}: Progettazione e codifica dei servizi di \textit{backend} (gestione utenti, autenticazione, storico chat, servizi di ricerca e server MCP) tramite il framework \textit{FastAPI}, strutturati secondo i dettami dell'architettura esagonale per disaccoppiare la logica di business dalle dipendenze tecnologiche.
    \item \textbf{Sviluppo di un'interfaccia utente reattiva}: Realizzazione del \textit{frontend} dell'applicazione mediante il \gls{frameworkg} \textit{React} (adottando il pattern MVVM), che includa una pagina per l'interazione in stile chatbot con il modello e la visualizzazione in tempo reale, tramite \gls{sseg}, dei singoli passaggi logici eseguiti dal \textit{crawler}.
    \item \textbf{Gestione della sicurezza e dell'autenticazione}: Implementazione di un sistema sicuro di controllo degli accessi basato su \gls{jwtg}, che regoli l'emissione, la validazione e la revoca di \textit{access token} e \textit{refresh token}, associando permessi specifici in base al ruolo dell'utente (es. utente standard ed amministratore).
    \item \textbf{Containerizzazione e orchestrazione dell'infrastruttura}: Utilizzo della piattaforma \textit{Docker} per l'isolamento di ciascun microservizio e del database relazionale, garantendo la parità degli ambienti, la portabilità del software e il rispetto dei vincoli di interoperabilità aziendali.
    \item \textbf{Verifica e documentazione}: Conduzione di \textit{test} di unità e di integrazione per validare il corretto funzionamento del \textit{routing} dei \textit{MCP tools} e stesura della documentazione tecnica finale.
\end{itemize}

\subsection{Obiettivi desiderabili}
Gli obiettivi desiderabili comprendono tutte quelle funzionalità aggiuntive, attività di ottimizzazione e strumenti di monitoraggio volti ad elevare la qualità del prodotto finale e la sua esperienza d'uso:

\begin{itemize}
    \item \textbf{Implementazione di un sistema di caching intelligente}: Sviluppo di un meccanismo di persistenza temporanea delle richieste per intercettare e gestire in modo efficiente i quesiti ridondanti formulati dagli utenti, abbattendo i tempi di latenza complessivi del sistema e riducendo il consumo di \textit{token} verso le \gls{apig} del modello linguistico.
    \item \textbf{Integrazione di sistemi di tracing avanzati}: Introduzione e configurazione della piattaforma \textit{LangSmith} all'interno del servizio \gls{mcpg}, allo scopo di monitorare, tracciare ed analizzare nel dettaglio il flusso delle chiamate e le performance dei modelli di intelligenza artificiale impiegati.
    \item \textbf{Funzionalità avanzate di gestione e persistenza}: Estensione delle capacità del sistema attraverso la creazione di un pannello amministrativo per la gestione delle utenze, la personalizzazione del profilo utente e la memorizzazione persistente su database dello storico delle conversazioni.
    \item \textbf{Esportazione locale dei dati generati}: Sviluppo di una funzionalità che consenta all'utente di salvare e scaricare la risposta generata dall'assistente virtuale in un file locale (in conformità con i requisiti di utilità aziendale), gestendo in modo robusto gli eventuali errori di scrittura del file.
    \item \textbf{Ottimizzazione dell'User Experience (UX)}: Introduzione di accorgimenti grafici e funzionali volti a migliorare l'usabilità dell'interfaccia web, tra cui il supporto del tema chiaro/scuro, elementi visivi iconografici immediati e opzioni di copia rapida negli appunti o di \textit{download} per i testi generati.
\end{itemize}